% ============================================================
%  cap05_espressioni.tex — Matemagica
%  Capitolo 5: Espressioni aritmetiche
% ============================================================

\chapter{Espressioni aritmetiche}

% --- Obiettivi di apprendimento ---
\section*{Obiettivi di apprendimento}
\begin{itemize}[leftmargin=1.5em]
    \item Comprendere il concetto di espressione aritmetica.
    \item Applicare correttamente le regole di precedenza tra le operazioni.
    \item Risolvere espressioni con parentesi tonde, quadre e graffe.
    \item Distinguere i due ruoli delle parentesi graffe: notazione insiemistica
          e raggruppamento nelle espressioni.
    \item Confrontare la notazione con sole parentesi tonde annidate e quella
          con i tre ordini di parentesi.
\end{itemize}

% ============================================================
\section{Che cos'è un'espressione aritmetica}
% ============================================================

\begin{definizione}
    Un'\textbf{espressione aritmetica} è una scrittura matematica che collega
    due o più numeri tramite segni di operazione. È una catena di operazioni il
    cui risultato finale è un unico numero, detto \textbf{valore dell'espressione}.
\end{definizione}

\begin{esempio}
    Le seguenti sono espressioni aritmetiche:
    \[
        12 + 5 \times 3 \qquad\qquad
        (8 - 2) \times 4 + 6 : 2 \qquad\qquad
        \{3 \times [10 - (2 + 1)]\} + 5
    \]
\end{esempio}

\begin{attenzione}
    \textbf{Un'espressione non è un'equazione.} Un'espressione aritmetica non
    contiene il segno di uguale: va \emph{calcolata}, non \emph{risolta}.
    Il segno $=$ appare solo nei passaggi intermedi che mostriamo durante il calcolo.
\end{attenzione}

% ============================================================
\section{Le regole di precedenza}
% ============================================================

Per ottenere un risultato univoco, ogni persona che calcola la stessa
espressione deve seguire le stesse regole. Queste si chiamano
\textbf{regole di precedenza}.

\begin{regola}
    \textbf{Ordine di esecuzione delle operazioni.}
    \begin{enumerate}[leftmargin=1.5em]
        \item Si risolvono prima le operazioni racchiuse fra \textbf{parentesi},
              cominciando da quelle più interne.
        \item Fra le operazioni senza parentesi, \textbf{moltiplicazione} e
              \textbf{divisione} hanno la precedenza su addizione e sottrazione.
        \item Fra operazioni dello stesso livello (solo addizioni e sottrazioni,
              oppure solo moltiplicazioni e divisioni) si procede
              \textbf{da sinistra verso destra}.
    \end{enumerate}
\end{regola}

\begin{esempio}
    Calcolare $12 + 5 \times 3 - 8 : 4$.

    Non ci sono parentesi, quindi si eseguono prima le moltiplicazioni e le divisioni,
    da sinistra verso destra, poi le addizioni e le sottrazioni:
    \begin{align*}
        12 + 5 \times 3 - 8 : 4
         & = 12 + 15 - 2 \\
         & = 27 - 2      \\
         & = 25
    \end{align*}
\end{esempio}

\begin{attenzione}
    \textbf{Errore comune — procedere semplicemente da sinistra verso destra.}
    Senza rispettare le precedenze, $12 + 5 \times 3$ potrebbe essere calcolato
    come $(12 + 5) \times 3 = 51$, risultato scorretto.
    La risposta corretta è $12 + 15 = 27$, perché la moltiplicazione precede
    l'addizione. \quad\textit{(corretto)}
\end{attenzione}

% ============================================================
\section{Le parentesi nelle espressioni}
% ============================================================

Le parentesi servono a modificare l'ordine naturale delle operazioni:
tutto ciò che è racchiuso fra una coppia di parentesi va calcolato
\emph{prima} del resto.

\subsection{Parentesi tonde annidate}

Il modo più semplice di indicare più livelli di raggruppamento è usare
più coppie di parentesi \textbf{tonde}, una dentro l'altra.
Questa scrittura si dice \textbf{annidata} (dal latino \emph{nidus}, nido:
ogni coppia di parentesi «contiene» quella più interna come un nido).

\begin{regola}
    Con le sole parentesi tonde annidate, si procede sempre dalla
    \textbf{parentesi più interna} verso quella più esterna.
\end{regola}

\begin{esempio}
    Calcolare $((3 + 2) \times (4 - 1) + 6) \times 2$.

    Si individua la parentesi più esterna e si lavora dall'interno:
    \begin{align*}
        ((3 + 2) \times (4 - 1) + 6) \times 2
         & = (5 \times 3 + 6) \times 2 &  & \text{[parentesi più interne]}     \\
         & = (15 + 6) \times 2         &  & \text{[moltiplicazione]}           \\
         & = 21 \times 2               &  & \text{[addizione nella parentesi]} \\
         & = 42
    \end{align*}
\end{esempio}

\begin{attenzione}
    \textbf{La notazione con sole tonde annidate è corretta ma poco leggibile.}
    Espressioni come $((8-(3+1))\times(5-(2+1)))+4$ sono difficili da leggere:
    è faticoso contare quante parentesi si aprono e quante si chiudono,
    e un errore di conteggio porta a un risultato sbagliato.
    Per questo motivo, nella pratica matematica si usa una notazione più chiara,
    basata su tre tipi di parentesi distinte.
\end{attenzione}

\subsection{I tre ordini di parentesi: tonde, quadre e graffe}

Per rendere le espressioni più leggibili e meno soggette a errori di lettura,
la matematica utilizza tre simboli distinti, usati dall'interno verso l'esterno:

\begin{center}
    \begin{tabular}{cll}
        \toprule
        \textbf{Simbolo} & \textbf{Nome}             & \textbf{Livello}    \\
        \midrule
        $(\quad)$        & Parentesi \textbf{tonde}  & Primo (più interno) \\
        $[\quad]$        & Parentesi \textbf{quadre} & Secondo             \\
        $\{\quad\}$      & Parentesi \textbf{graffe} & Terzo (più esterno) \\
        \bottomrule
    \end{tabular}
\end{center}

\begin{regola}
    Con i tre ordini di parentesi, si procede sempre dall'interno verso
    l'esterno, rispettando la gerarchia:
    \[
        \text{tonde} \longrightarrow \text{quadre} \longrightarrow \text{graffe}
    \]
    All'interno di ciascuna coppia di parentesi valgono le normali regole
    di precedenza: prima moltiplicazioni e divisioni, poi addizioni e sottrazioni.
\end{regola}

\begin{esempio}
    Calcolare $\{3 \times [10 - (2 + 1)] + 5\} : 4$.

    \textbf{Passo 1} — parentesi tonde (più interne):
    \[
        \{3 \times [10 - 3] + 5\} : 4
    \]
    \textbf{Passo 2} — parentesi quadre:
    \[
        \{3 \times 7 + 5\} : 4
        = \{21 + 5\} : 4
        = \{26\} : 4
    \]
    \textbf{Passo 3} — parentesi graffe:
    \[
        26 : 4 = 6{,}5
    \]
\end{esempio}

\begin{esempio}
    Calcolare $\{5 \times [(12 - 4) : (3 - 1)] - 2\} + 3 \times 4$.

    \textbf{Passo 1} — parentesi tonde (ce ne sono due, allo stesso livello):
    \begin{align*}
        \{5 \times [8 : 2] - 2\} + 3 \times 4
    \end{align*}
    \textbf{Passo 2} — parentesi quadre:
    \begin{align*}
        \{5 \times 4 - 2\} + 3 \times 4
        = \{20 - 2\} + 3 \times 4
        = \{18\} + 3 \times 4
    \end{align*}
    \textbf{Passo 3} — fuori da tutte le parentesi, moltiplicazione prima:
    \begin{align*}
        18 + 12 = 30
    \end{align*}
\end{esempio}

% ============================================================
\section{Confronto tra i due approcci}
% ============================================================

Come abbiamo visto, le due notazioni sono \textbf{equivalenti dal punto di
    vista matematico}: producono lo stesso risultato. La differenza è nella
leggibilità.

\begin{center}
    \begin{tabular}{p{0.42\linewidth} p{0.42\linewidth}}
        \toprule
        \textbf{Solo parentesi tonde annidate}                                    & \textbf{Tre ordini di parentesi} \\
        \midrule
        $((8-(3+1))\times(5-(2+1)))+4$                                            &
        $\{[8-(3+1)]\times[5-(2+1)]\}+4$                                                                             \\[4pt]
        Tutte le parentesi hanno la stessa forma: più livelli ci sono,
        più è difficile distinguerli a colpo d'occhio.                            &
        Ogni livello ha un simbolo diverso: è immediato vedere quale operazione
        va risolta prima.                                                                                            \\[4pt]
        Richiede di contare attentamente l'apertura e la chiusura di ogni coppia. &
        La forma visiva guida la lettura e riduce gli errori.                                                        \\
        \bottomrule
    \end{tabular}
\end{center}

\begin{sapevi}
    \textbf{Un doppio ruolo per le parentesi graffe.}

    Nel capitolo sugli insiemi abbiamo usato le parentesi graffe per elencare
    gli elementi di un insieme: $A = \{1;\; 3;\; 5\}$.
    Nelle espressioni aritmetiche, le stesse graffe svolgono un ruolo diverso:
    indicano il livello più esterno di raggruppamento.

    I due usi si distinguono dal contesto:
    \begin{itemize}[leftmargin=1.5em]
        \item Se le graffe contengono un elenco di numeri separati da punto e
              virgola, si tratta di un \textbf{insieme}.
        \item Se le graffe contengono un'operazione, si tratta di un
              \textbf{raggruppamento} in un'espressione.
    \end{itemize}
    In matematica è frequente che uno stesso simbolo abbia significati diversi
    a seconda del contesto: riconoscere il contesto è parte del ragionamento
    matematico.
\end{sapevi}

% ============================================================
\section{Le parentesi e la proprietà distributiva}
% ============================================================

Le parentesi non servono solo a stabilire l'ordine delle operazioni:
hanno anche un legame profondo con la \textbf{proprietà distributiva}
della moltiplicazione rispetto all'addizione e alla sottrazione.

\begin{regola}
    \textbf{Proprietà distributiva.}
    \[
        a \times (b + c) = a \times b + a \times c
        \qquad\qquad
        a \times (b - c) = a \times b - a \times c
    \]
    Questo significa che una moltiplicazione che «abbraccia» una somma
    (o una differenza) può essere distribuita su ciascun termine.
\end{regola}

\begin{esempio}
    Verificare che $4 \times (5 + 3) = 4 \times 5 + 4 \times 3$.
    \begin{align*}
        4 \times (5 + 3)        & = 4 \times 8 = 32                       \\
        4 \times 5 + 4 \times 3 & = 20 + 12 = 32 \quad\textit{(corretto)}
    \end{align*}
\end{esempio}

\begin{attenzione}
    \textbf{La proprietà distributiva NON vale per la divisione a sinistra.}
    \[
        (8 + 4) : 2 = 12 : 2 = 6
        \qquad\text{ma}\qquad
        8 : 2 + 4 : 2 = 4 + 2 = 6 \quad \text{(in questo caso funziona)}
    \]
    \[
        12 : (4 + 2) = 12 : 6 = 2
        \qquad\text{ma}\qquad
        12 : 4 + 12 : 2 = 3 + 6 = 9 \quad \text{(risultato diverso!)}
    \]
    Quando il \emph{divisore} contiene una somma racchiusa fra parentesi,
    la distribuzione non è applicabile.
\end{attenzione}

% ============================================================
\section*{Approfondimento — La convenzione internazionale}
% ============================================================

\begin{sapevi}
    \textbf{La posizione dell'American Physical Society.}

    L'\emph{American Physical Society} (APS), una delle principali organizzazioni
    scientifiche mondiali, ha codificato nelle proprie linee guida editoriali la
    convenzione che da secoli i matematici adottano per le parentesi nelle
    espressioni complesse.

    La regola stabilisce che, quando la forma delle parentesi non ha un significato
    speciale (come accade nelle espressioni aritmetiche), si deve lavorare
    dall'interno verso l'esterno seguendo la sequenza:
    \[
        \{ \;[\; (\quad) \;]\; \}
    \]
    e che \textbf{l'annidamento di sole parentesi tonde va evitato},
    perché rende la struttura dell'espressione difficile da leggere e da verificare.

    Questa non è una preferenza stilistica: è una convenzione che riduce
    gli errori e garantisce che la stessa espressione venga letta nello stesso
    modo da chiunque, indipendentemente dalla lingua o dalla scuola di appartenenza.
    Il fatto che alcuni testi scolastici introducano solo le parentesi tonde
    annidate è quindi una semplificazione didattica temporanea, non una regola
    matematica definitiva.
\end{sapevi}

% ============================================================
\section*{Riepilogo}
% ============================================================

\begin{regola}
    \textbf{Riepilogo del capitolo.}
    \begin{itemize}[leftmargin=1.5em]
        \item Un'\textbf{espressione aritmetica} è una catena di operazioni il
              cui calcolo segue regole precise.
        \item \textbf{Regole di precedenza}: prima le parentesi (dall'interno
              verso l'esterno), poi moltiplicazione e divisione, infine addizione
              e sottrazione; a parità di livello, da sinistra verso destra.
        \item Le parentesi si usano in tre ordini: \textbf{tonde} $(\,)$,
              \textbf{quadre} $[\,]$, \textbf{graffe} $\{\,\}$.
        \item Le due notazioni (sole tonde annidate / tre ordini) sono
              matematicamente equivalenti, ma la notazione a tre ordini è
              più leggibile e corrisponde alla convenzione internazionale.
        \item Le parentesi \textbf{graffe} hanno un doppio ruolo: indicano
              un insieme (con elementi separati da punto e virgola) oppure
              il livello più esterno di raggruppamento in un'espressione.
        \item La \textbf{proprietà distributiva} lega parentesi e
              moltiplicazione: $a \times (b + c) = a \times b + a \times c$.
    \end{itemize}
\end{regola}
